\documentclass[10pt,a4paper]{amsart}
\usepackage{amsmath,amssymb,amsthm,graphicx}
% Retained preamble for the Stacks Project chapter snapshot in this directory.
% The chapter's own preamble.tex is a copy of the shared Stacks Project preamble; the block
% below is sliced MECHANICALLY, line for line, from it. Nothing was added, removed or
% redefined; the wrapper class, the xr-hyper/hyperref wiring and the externaldocument list
% are omitted because the assembler supplies the document class.
%
% lines taken from preamble.tex: 3, 5, 9, 11, 12, 156, 157, 158, ... (51 lines total)
%
% package coverage restored from the chapter's own preamble (2026-10-13): xy, xy-2cell

\usepackage[all]{xy}
\xyoption{2cell}
\UseAllTwocells

\newenvironment{reference}{\comment}{\endcomment}
\newenvironment{slogan}{\comment}{\endcomment}
\newenvironment{history}{\comment}{\endcomment}
\theoremstyle{plain}
\newtheorem{theorem}[subsection]{Theorem}
\newtheorem{proposition}[subsection]{Proposition}
\newtheorem{lemma}[subsection]{Lemma}
\theoremstyle{definition}
\newtheorem{definition}[subsection]{Definition}
\newtheorem{example}[subsection]{Example}
\newtheorem{exercise}[subsection]{Exercise}
\newtheorem{situation}[subsection]{Situation}
\theoremstyle{remark}
\newtheorem{remark}[subsection]{Remark}
\newtheorem{remarks}[subsection]{Remarks}
\numberwithin{equation}{subsection}
\def\lim{\mathop{\mathrm{lim}}\nolimits}
\def\colim{\mathop{\mathrm{colim}}\nolimits}
\def\Spec{\mathop{\mathrm{Spec}}}
\def\Hom{\mathop{\mathrm{Hom}}\nolimits}
\def\Ext{\mathop{\mathrm{Ext}}\nolimits}
\def\SheafHom{\mathop{\mathcal{H}\!\mathit{om}}\nolimits}
\def\SheafExt{\mathop{\mathcal{E}\!\mathit{xt}}\nolimits}
\def\Sch{\mathit{Sch}}
\def\Mor{\mathop{\mathrm{Mor}}\nolimits}
\def\Ob{\mathop{\mathrm{Ob}}\nolimits}
\def\Sh{\mathop{\mathit{Sh}}\nolimits}
\def\NL{\mathop{N\!L}\nolimits}
\def\CH{\mathop{\mathrm{CH}}\nolimits}
\def\proetale{{pro\text{-}\acute{e}tale}}
\def\etale{{\acute{e}tale}}
\def\QCoh{\mathit{QCoh}}
\def\Ker{\mathop{\mathrm{Ker}}}
\def\Im{\mathop{\mathrm{Im}}}
\def\Coker{\mathop{\mathrm{Coker}}}
\def\Coim{\mathop{\mathrm{Coim}}}
\DeclareMathSymbol{\boxtimes}{\mathbin}{AMSa}{"02}
\def\QCohstack{\mathcal{QC}\!\mathit{oh}}
\def\Cohstack{\mathcal{C}\!\mathit{oh}}
\def\Spacesstack{\mathcal{S}\!\mathit{paces}}
\def\Quotfunctor{\mathrm{Quot}}
\def\Hilbfunctor{\mathrm{Hilb}}
\def\Curvesstack{\mathcal{C}\!\mathit{urves}}
\def\Polarizedstack{\mathcal{P}\!\mathit{olarized}}
\def\Complexesstack{\mathcal{C}\!\mathit{omplexes}}
\def\Pic{\mathop{\mathrm{Pic}}\nolimits}
\def\Picardstack{\mathcal{P}\!\mathit{ic}}
\def\Picardfunctor{\mathrm{Pic}}
\def\Deformationcategory{\mathcal{D}\!\mathit{ef}}
\usepackage{cleveref}
\crefname{originallemma}{original lemma}{original lemmas}
\Crefname{originallemma}{Original lemma}{Original lemmas}
\makeatletter
\newlabel{defos-lemma-choices-ringed-spaces}{{08UC}{source}{}{sourcepointer-1}{}}
\newlabel{defos-lemma-choices-ringed-spaces@cref}{{[originallemma][0][]08UC}{[1][0][]source}{}{}{}}
\makeatother
\crefname{originallemma}{original lemma}{original lemmas}
\Crefname{originallemma}{Original lemma}{Original lemmas}
\makeatletter
\newlabel{defos-lemma-huge-diagram-ringed-spaces}{{08U8}{source}{}{sourcepointer-2}{}}
\newlabel{defos-lemma-huge-diagram-ringed-spaces@cref}{{[originallemma][0][]08U8}{[1][0][]source}{}{}{}}
\makeatother
\crefname{originallemma}{original lemma}{original lemmas}
\Crefname{originallemma}{Original lemma}{Original lemmas}
\makeatletter
\newlabel{defos-lemma-parametrize-solutions-ringed-spaces}{{0GQ3}{source}{}{sourcepointer-3}{}}
\newlabel{defos-lemma-parametrize-solutions-ringed-spaces@cref}{{[originallemma][0][]0GQ3}{[1][0][]source}{}{}{}}
\makeatother
\crefname{originalremark}{original remark}{original remarks}
\Crefname{originalremark}{Original remark}{Original remarks}
\makeatletter
\newlabel{defos-remark-parametrize-solutions-ringed-spaces}{{0GQ4}{source}{}{sourcepointer-4}{}}
\newlabel{defos-remark-parametrize-solutions-ringed-spaces@cref}{{[originalremark][0][]0GQ4}{[1][0][]source}{}{}{}}
\makeatother
\crefname{originalequation}{original equation}{original equations}
\Crefname{originalequation}{Original equation}{Original equations}
\makeatletter
\newlabel{equation-to-solve-ringed-spaces}{{08UY}{source}{}{sourcepointer-5}{}}
\newlabel{equation-to-solve-ringed-spaces@cref}{{[originalequation][0][]08UY}{[1][0][]source}{}{}{}}
\makeatother
\crefname{originallemma}{original lemma}{original lemmas}
\Crefname{originallemma}{Original lemma}{Original lemmas}
\makeatletter
\newlabel{lemma-H0-L-morphism-ringed-spaces}{{08UV}{source}{}{sourcepointer-6}{}}
\newlabel{lemma-H0-L-morphism-ringed-spaces@cref}{{[originallemma][0][]08UV}{[1][0][]source}{}{}{}}
\makeatother
\crefname{originallemma}{original lemma}{original lemmas}
\Crefname{originallemma}{Original lemma}{Original lemmas}
\makeatletter
\newlabel{lemma-compare-cotangent-complex-with-naive-ringed-spaces}{{08UW}{source}{}{sourcepointer-7}{}}
\newlabel{lemma-compare-cotangent-complex-with-naive-ringed-spaces@cref}{{[originallemma][0][]08UW}{[1][0][]source}{}{}{}}
\makeatother
\crefname{originallemma}{original lemma}{original lemmas}
\Crefname{originallemma}{Original lemma}{Original lemmas}
\makeatletter
\newlabel{lemma-triangle-ringed-spaces}{{08T4}{source}{}{sourcepointer-8}{}}
\newlabel{lemma-triangle-ringed-spaces@cref}{{[originallemma][0][]08T4}{[1][0][]source}{}{}{}}
\makeatother
\title{The obstruction class for deforming a morphism of ringed spaces over a thickening}
\author{The Stacks Project Authors}
\date{stacks-project:cotangent.tex:local-snapshot-20261012; GNU FDL 1.2 or later}
\begin{document}
\maketitle
\setcounter{subsection}{1}
\noindent\textit{Editorial note.} All delivered mathematical text is the author's own, copied from the pinned Stacks Project chapter file cotangent.tex; no statement or proof step was rewritten, reordered or removed.\par\medskip
\noindent\textit{Editorial note.} The statement refers to the situation and the equation-to-solve fixed in the preceding section; the proof recalls the two identifications it uses.\par\medskip
\noindent\textit{Editorial note.} The eight references of the counted spans are recorded as non-load-bearing navigation pointers; the argument that consumes them stays inside the counted proof span.\par\medskip
\noindent\textit{Editorial note.} Preamble: the retained preamble of this chapter (sources/stacks-cotangent/retained-preamble.tex, lines 1-63), containing the Xy-pic setup, the comment environments, the theorem declarations and the macros including \textbackslash{}NL.\par\medskip
\noindent\textit{Allowed standard prerequisites (editorial).}\begin{itemize}
\item The identification of the naive cotangent complex with the truncation of the cotangent complex for ringed spaces. Lemma lemma-compare-cotangent-complex-with-naive-ringed-spaces, used at proof lines 3854-3856. Transfers the deformation-theoretic statements.
\item The identification H\textasciicircum{}0(L\_\{X/S\}) = Omega\_\{X/S\}. Lemma lemma-H0-L-morphism-ringed-spaces, used at proof lines 3857-3859. Is the second identification used for parts (2) and (3).
\item The deformation-theory statements for morphisms of ringed spaces. Deformation Theory, Lemmas defos-lemma-parametrize-solutions-ringed-spaces, defos-lemma-huge-diagram-ringed-spaces and defos-lemma-choices-ringed-spaces, used at proof lines 3859-3861 and 3864-3870. Supply parts (2), (3) and the class xi'.
\item The distinguished triangle for X -> S -> S'. Lemma lemma-triangle-ringed-spaces, used at proof lines 3884-3891. Gives the long exact sequence producing the obstruction.
\end{itemize}\medskip
\ifdefined\hypertarget\hypertarget{sourcepointer-1}{}\fi
\noindent\textit{Source reference.} Original lemma 08UC: focused-seeds/AGP/sources/stacks-snapshot-20261011/tags.txt, source lines 11195--11195. Navigation pointer into the Stacks Project; the standard fact it supplies is applied inside the counted proof span and the argument resting on it stays there.\par\medskip
\ifdefined\hypertarget\hypertarget{sourcepointer-2}{}\fi
\noindent\textit{Source reference.} Original lemma 08U8: focused-seeds/AGP/sources/stacks-snapshot-20261011/tags.txt, source lines 11191--11191. Navigation pointer into the Stacks Project; the standard fact it supplies is applied inside the counted proof span and the argument resting on it stays there.\par\medskip
\ifdefined\hypertarget\hypertarget{sourcepointer-3}{}\fi
\noindent\textit{Source reference.} Original lemma 0GQ3: focused-seeds/AGP/sources/stacks-snapshot-20261011/tags.txt, source lines 20848--20848. Navigation pointer into the Stacks Project; the standard fact it supplies is applied inside the counted proof span and the argument resting on it stays there.\par\medskip
\ifdefined\hypertarget\hypertarget{sourcepointer-4}{}\fi
\noindent\textit{Source reference.} Original remark 0GQ4: focused-seeds/AGP/sources/stacks-snapshot-20261011/tags.txt, source lines 20849--20849. Navigation pointer into the Stacks Project; the standard fact it supplies is applied inside the counted proof span and the argument resting on it stays there.\par\medskip
\ifdefined\hypertarget\hypertarget{sourcepointer-5}{}\fi
\noindent\textit{Source reference.} Original equation 08UY: focused-seeds/AGP/sources/stacks-snapshot-20261011/tags.txt, source lines 11216--11216. Navigation pointer into the Stacks Project; the standard fact it supplies is applied inside the counted proof span and the argument resting on it stays there.\par\medskip
\ifdefined\hypertarget\hypertarget{sourcepointer-6}{}\fi
\noindent\textit{Source reference.} Original lemma 08UV: focused-seeds/AGP/sources/stacks-snapshot-20261011/tags.txt, source lines 11213--11213. Navigation pointer into the Stacks Project; the standard fact it supplies is applied inside the counted proof span and the argument resting on it stays there.\par\medskip
\ifdefined\hypertarget\hypertarget{sourcepointer-7}{}\fi
\noindent\textit{Source reference.} Original lemma 08UW: focused-seeds/AGP/sources/stacks-snapshot-20261011/tags.txt, source lines 11214--11214. Navigation pointer into the Stacks Project; the standard fact it supplies is applied inside the counted proof span and the argument resting on it stays there.\par\medskip
\ifdefined\hypertarget\hypertarget{sourcepointer-8}{}\fi
\noindent\textit{Source reference.} Original lemma 08T4: focused-seeds/AGP/sources/stacks-snapshot-20261011/tags.txt, source lines 11152--11152. Navigation pointer into the Stacks Project; the standard fact it supplies is applied inside the counted proof span and the argument resting on it stays there.\par\medskip
\begin{lemma}
\label{lemma-find-obstruction-ringed-spaces}
In the situation above we have
\begin{enumerate}
\item There is a canonical element
$\xi \in \Ext^2_{\mathcal{O}_X}(L_{X/S}, \mathcal{G})$
whose vanishing is a sufficient and necessary condition for the existence
of a solution to (\ref{equation-to-solve-ringed-spaces}).
\item If there exists a solution, then the set of
isomorphism classes of solutions is principal homogeneous under
$\Ext^1_{\mathcal{O}_X}(L_{X/S}, \mathcal{G})$.
\item Given a solution $X'$, the set of automorphisms of $X'$
fitting into (\ref{equation-to-solve-ringed-spaces}) is canonically isomorphic
to $\Ext^0_{\mathcal{O}_X}(L_{X/S}, \mathcal{G})$.
\end{enumerate}
\end{lemma}

\begin{proof}
Via the identifications $\NL_{X/S} = \tau_{\geq -1}L_{X/S}$
(Lemma \ref{lemma-compare-cotangent-complex-with-naive-ringed-spaces})
and
$H^0(L_{X/S}) = \Omega_{X/S}$
(Lemma \ref{lemma-H0-L-morphism-ringed-spaces})
we have seen parts (2) and (3) in
Deformation Theory, Lemmas \ref{defos-lemma-huge-diagram-ringed-spaces} and
\ref{defos-lemma-choices-ringed-spaces}.

\medskip\noindent
Proof of (1). Roughly speaking, this follows from the discussion in
Deformation Theory, Remark
\ref{defos-remark-parametrize-solutions-ringed-spaces}
by replacing the naive cotangent complex by the full cotangent complex.
Here is a more detailed explanation. By
Deformation Theory, Lemma \ref{defos-lemma-parametrize-solutions-ringed-spaces}
there exists an element
$$
\xi' \in
\Ext^1_{\mathcal{O}_X}(Lf^*\NL_{S/S'}, \mathcal{G}) =
\Ext^1_{\mathcal{O}_X}(Lf^*L_{S/S'}, \mathcal{G})
$$
such that a solution exists if and only if this element is in
the image of the map
$$
\Ext^1_{\mathcal{O}_X}(NL_{X/S'}, \mathcal{G}) =
\Ext^1_{\mathcal{O}_X}(L_{X/S'}, \mathcal{G})
\longrightarrow
\Ext^1_{\mathcal{O}_X}(Lf^*L_{S/S'}, \mathcal{G})
$$
The distinguished triangle of Lemma \ref{lemma-triangle-ringed-spaces}
for $X \to S \to S'$ gives rise to a long exact sequence
$$
\ldots \to
\Ext^1_{\mathcal{O}_X}(L_{X/S'}, \mathcal{G}) \to
\Ext^1_{\mathcal{O}_X}(Lf^*L_{S/S'}, \mathcal{G}) \to
\Ext^2_{\mathcal{O}_X}(L_{X/S}, \mathcal{G}) \to \ldots
$$
Hence taking $\xi$ the image of $\xi'$ works.
\end{proof}

\end{document}
